\documentclass[11pt]{article}
\usepackage[margin=1in]{geometry}
\usepackage{amsmath,amssymb,amsthm}
\usepackage{enumitem}

\newcommand{\Inf}{\operatorname{Inf}}
\newcommand{\Maj}{\operatorname{Maj}}
\newcommand{\E}{\mathbf E}
\newcommand{\I}{\mathbf I}
\newcommand{\Wone}{\mathbf W^{1}}

\title{Solutions to O'Donnell, Analysis of Boolean Functions, Exercise 2.22}
\author{}
\date{}

\begin{document}
\maketitle

Throughout, for odd $n$ we write $n=2m+1$ and use the
standard majority function
\[
  \Maj_n(x)=\operatorname{sgn}(x_1+\cdots+x_n),
  \qquad x\in\{-1,1\}^n.
\]
Thus there is no tie when $n$ is odd.

\begin{enumerate}[label=\textbf{(\alph*)}]
\item
We prove
\[
  \Inf_i[\Maj_n]
    =
  \binom{n-1}{(n-1)/2}2^{1-n}
  \qquad \text{for every } i\in[n].
\]
Fix $i$ and condition on all coordinates except $x_i$.  Among the
remaining $n-1=2m$ coordinates, let $k$ be the number of $+1$'s.

If $k>m$, then the other coordinates already have a strict positive
majority, so changing $x_i$ cannot change the value of $\Maj_n$.  If
$k<m$, the other coordinates already have a strict negative majority,
and again changing $x_i$ cannot change the value of $\Maj_n$.  The only
case in which $i$ is pivotal is $k=m$.

When $k=m$, the other $2m$ coordinates are tied.  Then setting $x_i=1$
makes the total sum equal to $1$, while setting $x_i=-1$ makes the total
sum equal to $-1$.  Thus the value of majority changes.  Hence $i$ is
pivotal exactly when the remaining $2m$ bits contain exactly $m$ many
$+1$'s.

There are $\binom{2m}{m}$ assignments of the remaining coordinates with
this property, and $2^{2m}$ total assignments of the remaining
coordinates.  Therefore
\[
  \Inf_i[\Maj_{2m+1}]
    =
  \frac{\binom{2m}{m}}{2^{2m}}
    =
  \binom{n-1}{(n-1)/2}2^{1-n}.
\]

\item
Let
\[
  a_m=\Inf_1[\Maj_{2m+1}]
      =\frac{\binom{2m}{m}}{4^m}.
\]
Then
\[
  \frac{a_{m+1}}{a_m}
   =
  \frac{\binom{2m+2}{m+1}}{4\binom{2m}{m}}
   =
  \frac{(2m+2)(2m+1)}{4(m+1)^2}
   =
  \frac{2m+1}{2m+2}
   <1.
\]
Thus $a_m$ is strictly decreasing in $m$, equivalently
$\Inf_1[\Maj_n]$ is decreasing as a function of odd $n$.

\item
We again use
\[
  \Inf_1[\Maj_{2m+1}]
    =
  \frac{\binom{2m}{m}}{4^m}
    =
  \frac{(2m)!}{(m!)^2 4^m}.
\]
Stirling's formula gives
\[
  m! = \sqrt{2\pi m}\,(m/e)^m(1+O(1/m)).
\]
Applying this also to $(2m)!$,
\[
\begin{aligned}
  \binom{2m}{m}
  &=
  \frac{\sqrt{4\pi m}\,(2m/e)^{2m}(1+O(1/m))}
       {2\pi m\,(m/e)^{2m}(1+O(1/m))}  \\
  &=
  \frac{4^m}{\sqrt{\pi m}}\,(1+O(1/m)).
\end{aligned}
\]
Hence
\[
  \Inf_1[\Maj_{2m+1}]
    =
  \frac{1}{\sqrt{\pi m}}+O(m^{-3/2}).
\]
Since $m=(n-1)/2$,
\[
  \frac{1}{\sqrt{\pi m}}
    =
  \sqrt{\frac{2}{\pi}}\frac{1}{\sqrt{n-1}}
    =
  \sqrt{\frac{2}{\pi}}\frac{1}{\sqrt n}
      +O(n^{-3/2}).
\]
Therefore
\[
  \Inf_1[\Maj_n]
    =
  \sqrt{\frac{2}{\pi}}\frac{1}{\sqrt n}
    +O(n^{-3/2})
\]
for odd $n$.

\item
For monotone Boolean functions, Proposition 2.21 gives
$\widehat f(i)=\Inf_i[f]$.  Majority is monotone and symmetric, so all
degree-one Fourier coefficients are equal.  Thus
\[
  \Wone[\Maj_{2m+1}]
    =
  \sum_{i=1}^{2m+1}\widehat{\Maj_{2m+1}}(i)^2
    =
  (2m+1)a_m^2,
  \qquad
  a_m=\frac{\binom{2m}{m}}{4^m}.
\]
Let $B_m=(2m+1)a_m^2$.  Then
\[
  \frac{B_{m+1}}{B_m}
    =
  \frac{2m+3}{2m+1}
  \left(\frac{a_{m+1}}{a_m}\right)^2
    =
  \frac{2m+3}{2m+1}
  \left(\frac{2m+1}{2m+2}\right)^2
    =
  1-\frac{1}{(2m+2)^2}
    <1.
\]
Thus $B_m$ decreases with $m$.  Part (c) implies
\[
  B_m
    =
  (2m+1)
  \left(
    \sqrt{\frac{2}{\pi}}\frac{1}{\sqrt{2m+1}}
    +O((2m+1)^{-3/2})
  \right)^2
    =
  \frac{2}{\pi}+O(n^{-1}).
\]
In particular $B_m\to 2/\pi$.  Since $B_m$ is decreasing to this limit,
\[
  \frac{2}{\pi}
    \le
  \Wone[\Maj_n]
    \le
  \frac{2}{\pi}+O(n^{-1})
\]
for odd $n$.

\item
The total influence is the sum of the individual influences:
\[
  \I[\Maj_n]
    =
  \sum_{i=1}^n \Inf_i[\Maj_n]
    =
  n\,\Inf_1[\Maj_n].
\]
Using part (c),
\[
  \I[\Maj_n]
    =
  n\left(
    \sqrt{\frac{2}{\pi}}\frac{1}{\sqrt n}
    +O(n^{-3/2})
  \right)
    =
  \sqrt{\frac{2}{\pi}}\sqrt n+O(n^{-1/2}).
\]
For the lower bound, part (d) gives
\[
  \Wone[\Maj_n]
    =
  n\,\Inf_1[\Maj_n]^2
    \ge
  \frac{2}{\pi}.
\]
Since $\Inf_1[\Maj_n]\ge 0$,
\[
  \I[\Maj_n]
    =
  n\,\Inf_1[\Maj_n]
    =
  \sqrt{n\,\Wone[\Maj_n]}
    \ge
  \sqrt{\frac{2}{\pi}}\sqrt n.
\]
Combining the lower bound with the asymptotic upper estimate yields
\[
  \sqrt{\frac{2}{\pi}}\sqrt n
    \le
  \I[\Maj_n]
    \le
  \sqrt{\frac{2}{\pi}}\sqrt n+O(n^{-1/2})
\]
for odd $n$.

\item
Let $n=2m$ be even, and let
$f:\{-1,1\}^{2m}\to\{-1,1\}$ be any majority function: it agrees with
the sign of $x_1+\cdots+x_{2m}$ when this sum is nonzero, and it may
break ties arbitrarily when the sum is zero.

We count boundary edges.  A boundary edge is an unordered edge of the
hypercube whose endpoints have different $f$-values.  Total influence is
\[
  \I[f]=\frac{2\cdot \#\{\text{boundary edges of } f\}}{2^{2m}},
\]
because every boundary edge contributes once from each endpoint to the
average sensitivity.

If a vertex is not on the tie layer, then an edge can cross the boundary
only if its other endpoint is on the tie layer.  Thus every boundary edge
has a unique endpoint $x$ with
\[
  x_1+\cdots+x_{2m}=0.
\]
There are $\binom{2m}{m}$ such tie vertices.

Fix a tie vertex $x$.  It has exactly $m$ coordinates equal to $+1$ and
$m$ coordinates equal to $-1$.  If $f(x)=1$, then flipping one of the
$+1$ coordinates moves to total sum $-2$, where every majority function
has value $-1$; these are exactly $m$ boundary edges incident to $x$.
Flipping a $-1$ coordinate moves to total sum $+2$, where the value is
$+1$, so those edges are not boundary edges.  If $f(x)=-1$, the same
argument with signs reversed again gives exactly $m$ boundary edges
incident to $x$.

Hence every tie vertex contributes exactly $m$ boundary edges, independent
of the tie-breaking rule.  Therefore
\[
  \#\{\text{boundary edges}\}=m\binom{2m}{m},
\]
and
\[
  \I[f]
    =
  \frac{2m\binom{2m}{m}}{2^{2m}}.
\]
On the other hand, by part (a),
\[
  \I[\Maj_{2m-1}]
    =
  (2m-1)\frac{\binom{2m-2}{m-1}}{2^{2m-2}}.
\]
Using
\[
  \binom{2m}{m}
    =
  \frac{(2m)(2m-1)}{m^2}\binom{2m-2}{m-1},
\]
we get
\[
  \frac{2m\binom{2m}{m}}{2^{2m}}
    =
  \frac{(2m-1)\binom{2m-2}{m-1}}{2^{2m-2}}
    =
  \I[\Maj_{2m-1}].
\]
Thus
\[
  \I[f]=\I[\Maj_{n-1}].
\]
Finally, applying part (e) to the odd integer $n-1$ gives
\[
  \I[f]
    =
  \sqrt{\frac{2}{\pi}}\sqrt{n-1}+O(n^{-1/2})
    =
  \sqrt{\frac{2}{\pi}}\sqrt n+O(n^{-1/2}).
\]
\end{enumerate}

\end{document}
